\documentclass[../main.tex]{subfiles}
\begin{document}
%==========================================TIÊU ĐỀ TẬP========================================
\begin{table}[h]
  \begin{adjustwidth}{-1cm}{}
    \begin{tabular}{>{\centering\arraybackslash}p{6cm}|>{\raggedright\arraybackslash}p{12.5cm}}
      \multirow{5}{6cm}
      % Cột bên trái: ÉPISODE và số tập
      {\\[-40pt]
        \flaregothic\fontsize{20pt}{20pt}\selectfont \centering
        \color{eptype}ÉPISODE\color{black}\\
        \burbank\fontsize{72pt}{72pt}\selectfont
        \color{epnum}124\color{black}\\[6pt]
        \cafeteria\fontsize{20pt}{20pt}\selectfont\color{epnum}(10-4)\color{black}
        \\[18pt]
        % \flaregothic\fontsize{14pt}{14pt}\selectfont
        % \color{epattr}BIENVENUE, \uline{\color{Banpire} Mel} !
        % \flaregothic\fontsize{18pt}{18pt}\selectfont
        % \color{epattr}ÉPISODE PERDU
        \color{black}
      }
       &
      % Dòng thứ nhất: tên tập tiếng Nhật
      {\fotskip\fontsize{18pt}{18pt}\selectfont \ruby{見}{み}るの\ruby{禁}{きん}\ruby{止}{し}！
      }                                     \\[6pt]
      % Dòng thứ hai: tên tập tiếng Anh
       & {\cafeteria\fontsize{18pt}{18pt}\selectfont TV Dinner}                                                          \\[4pt]
      % Dòng thứ ba: tên tập tiếng Việt
       & {\vnmsans\fontsize{18pt}{18pt}\selectfont Chương trình thiếu nhi}                                               \\[8pt]
      % Dòng thứ tư: Nhân vật xuất hiện
       & {\sen{\fontsize{14pt}{14pt}\selectfont Track used: Sci-Tech (\ruby{科}{か}\ruby{学}{がく}\ruby{技}{ぎ}\ruby{術}{じゅつ})}}
      \\[6pt]
      % Dòng cuối cùng: Thời gian diễn ra của tập (thường sẽ là ngày phát hành)
       & {\continuum\fontsize{16pt}{16pt}\selectfont Ngày phát hành: 26/09/2021}                                         \\[18pt]
    \end{tabular}
  \end{adjustwidth}
\end{table}
%=========================================TRUYỆN CHÍNH========================================
\color{black}

Cuối tháng Chín, tiết trời mùa thu đã bắt đầu chạm ngõ rõ rệt. Những cơn gió mát lành lùa qua ô cửa, thổi lay nhẹ tấm rèm màu kem trong văn phòng. Dù nắng ban trưa vẫn còn vương chút oi ả cuối hạ, nhưng cái se se của buổi sáng và chiều đã nhắc nhở mọi người rằng mùa lá vàng rơi đang bắt đầu gõ cửa tới mọi ngóc ngách của thành phố.

Trong văn phòng - cũng là nơi sinh hoạt chung, Rushia và Watame đang ngồi trước chiếc TV cỡ lớn, mắt dán chặt vào màn hình, gương mặt chăm chú như thể đang xem một bộ phim điều tra phá án gay cấn, nhưng thật ra chỉ là chương trình giải trí thiếu nhi, đầy tiếng nhạc lục lạc và hoạt hình.

Ở phía đối diện, Choco đang ngồi thoải mái trên ghế bành, đọc một cuốn manga shoujo tình cảm mà cô vừa mượn từ Subaru. Kuon thì ngồi sát bên, laptop trên đùi, ngón tay lướt qua các cột số liệu kế toán nội bộ, ánh mắt tỉnh táo và tập trung.

Đột nhiên, ánh mắt của nàng Y tá liếc qua hai cô gái ngồi quá gần chiếc TV. Cô nhăn mặt. ``Hai đứa ngồi gần cái TV quá đấy. Lùi ra xa chút đi.''

Giọng nói dịu dàng nhưng rõ ràng không chấp nhận phản kháng. Rushia và Watame đồng thanh như hai học sinh tiểu học: ``Vâng\~{}!!''

Choco mỉm cười nhẹ, hài lòng quay lại với cuốn truyện. Kuon cũng liếc nhìn hai người, nhún vai một cái, rồi lại tiếp tục với bảng tính.

Chưa đầy năm phút sau, cô ngẩng đầu lần nữa, và chẳng ngạc nhiên mấy khi thấy hai ``thủ phạm'' vẫn ngồi y nguyên chỗ cũ, thậm chí còn rướn người gần hơn một chút. Cậu Tổng cũng phát hiện ra điều đó, lắc đầu.

\img{10-4a}

``Cô vừa mới bảo hai em lùi ra xa đi mà?'' nàng Y tá lên tiếng, giọng bực hơn.

``Nghe lời cô đi mấy đứa ơi,'' cậu chêm thêm.

``Ngay bây giờ,'' cô nhấn mạnh, tay gập sách lại như chuẩn bị dùng để ném.

Hai người trước màn hình, vẫn không rời mắt khỏi chương trình, lại đồng thanh lần nữa: ``Vâng\~{}!''

Choco thở dài, gập cuốn truyện lại, nhìn hai cô nàng với đôi mắt cá ươn: ``Chúng lại không nghe lời em rồi\dots''

Kuon uể oải gõ thêm vài dòng lệnh rồi nhìn sang: ``Anh chả biết chúng đang xem cái gì mà say sưa thế.''

``Để em thử cách này.''

Nàng Y tá đi tới chỗ của nàng Cừu và nàng Pháp sư, rướn người hỏi qua: ``Hai đứa muốn ăn gì nào?''

Mắt Rushia và Watame sáng rực như đèn pin. Trước khi họ kịp trả lời, Kuon để ý thấy một chiếc lục lạc đang lăn lóc trên bàn, và lập tức suy luận: ``\textit{Không lẽ là\dots\ chương trình đó?}''

Chưa đầy một giây sau, hai người đồng thanh: ``Vâng\~{}!''

*CỐC!* *CỐC!*

Hai cú gõ chính xác như dùng kỹ thuật karatedo rơi ngay vào đầu của hai ``nạn nhân.''

\img{10-4b}

``\textbf{Tập trung vào cô đi!}'' Choco quát nhẹ.

Từ phía xa, nhìn Watame và Rushia xoa đầu, thỉnh thoảng rít lên vì đau, Kuon cười khì: ``\textit{Xứng đáng bị uýnh lắm. Thế là còn nhẹ đấy.}''

Choco chỉ thẳng vào mắt mình để minh họa, mặt tỏ vẻ đáng sợ: ``Các em mà cứ say sưa thế này là hỏng mắt đấy.''

``Đúng đó,'' Kuon tiếp lời, vẻ mặt vừa nghiêm túc vừa mỉa mai, ``Cẩn thận bị cận nặng luôn chứ chẳng đùa.''

Bỗng nhiên, Rushia bật dậy, giơ một tay như đang biện hộ tại tòa án: ``Em xin phản đối, thưa Quý tòa!''

``\dots Gì cơ?'' cậu tròn mắt.

``Bọn em chỉ đang xem TV thôi, chứ bọn em không có say sưa gì đâu ạ!''

``Có khác quái gì nhau đâu em\dots''

``Đơn phản bác không được chấp thuận. Đừng có lý sự nữa, trật tự,'' Choco ra phán quyết, tay vẫy như thẩm phán xử án.

``\textit{Mấy người làm cứ như thể đang trong tòa xét xử thật ấy\dots}'' cậu lẩm bẩm, không biết nên phản ứng ra sao.

Không còn gì để phản bác, Cô đồng tóc xanh cắn nhẹ răng, nghiến nghiến như một cái máy nghiền giấy. Âm thanh kẽo kẹt vang lên khiến cả ba người còn lại rùng mình.

``Phản đối không được lại quay sang biểu tình hả? Không có đâu em ơi. Bọn anh dân chủ lắm,'' cậu nhún vai, cười gượng với cách mà nàng Pháp sư nhõng nhẽo.

Watame bên cạnh lặng lẽ nhìn đàn chị của mình đang nhai không khí, ánh mắt đầy ngưỡng mộ: ``Trông có vẻ vui đấy\dots''

``Và đừng nghĩ đến chuyện làm theo  em ấy,'' cậu quay sang nàng Cừu, lên giọng, ``Người ta sẽ nghĩ em bị điên đấy.''

``Vâng\~{} Dù gì em cũng không có lỗi gì mà\~{}'' nàng Cừu hát bài ca ngây thơ quen thuộc của mình.

Kuon chống cằm nhìn hai cô gái như thể đang nuôi trong nhà hai phiên bản hoạt hình nhân hóa của\dots\ tủ lạnh và máy giặt, và thở dài lần thứ ba trong mười phút.

\begin{center}
  \uline{Ngày hôm sau.}
\end{center}

Nắng thu rọi xuyên qua các tán cây bên hông trụ sở \hll, chiếu nhẹ lên mặt đường lát đá, phản chiếu lại những vệt sáng ấm áp. Kuon và Choco đang đi bộ trên vỉa hè hướng vào văn phòng chính, vừa trò chuyện vừa lật lại vài ghi chú liên quan đến lịch phát trực tiếp sắp tới.

``\dots Em cứ thử thêm mấy câu hỏi lịch sử, địa lý nữa là ổn rồi,'' cậu Tổng vừa đi vừa viết ghi chú vào ứng dụng lịch trên điện thoại.

``Vâng, vậy để em suy xét thêm\dots'' nàng Y tá đáp, tay đưa lên môi như đang cân nhắc sâu xa.

Cô vốn rất thích dạng stream đối đáp nhanh, những kiểu đố mẹo, vặn vẹo logic, hoặc chỉ đơn giản là kiểm tra kiến thức cơ bản của mọi người trong công ty. Và đoán xem ai cũng thích dạng này ngoài Kuon nào? Chính cậu là người đã ra đề Toán dịp Tết Nguyên Đán năm ngoái\footnote{Xem lại {\sen Partie 6}, {\flaregothic PERSONNALISÉ \vnmdisp MỘT}, quyển T.}, với mức độ được mọi người mô tả là ``bất khả thi đến mức đòi kiện,'' thậm chí đến cô em gái Ryuuhou của cậu cũng phải khóc thét dẫu cho việc đề của cô bé còn dễ hơn so với những người khác.

Cũng nhờ thế mà hai người có chung một sở thích không tưởng - giữa một nàng Quỷ Y tá khiêu gợi và một chàng quản lý bề ngoài tốt bụng nhưng cũng máu lạnh không kém.

Cả hai bước vào hành lang tầng ba và đồng thanh cất tiếng: ``Chào buổi sáng.''

Chưa kịp đặt cặp xuống bàn, Subaru từ trong văn phòng lao ra như tên bắn, mồ hôi nhễ nhại: ``\textbf{Nguy rồi, Choco-sen, Kuon-senpai!!}''

``Bình tĩnh lại đi chứ!'' Choco đưa tay giữ lấy vai nàng Vịt. Trong lúc đó, Kuon nhướng này: ``Sao thế, Suba-chan?''

``\textbf{Rushia và Watame đang ôm lấy cái TV kìa!!}'' Subaru xanh mặt, tay chỉ về bên trong văn phòng.

``\dots Lại nữa à?'' hai người lớn lập tức đổi sắc mặt.

Nói rồi, cả ba vội vã chạy đến văn phòng. Nhưng cảnh tượng trước mắt khiến họ phải đứng hình trong vài giây. Đúng là Rushia và Watame đang ghim chặt vào chiếc vô tuyến\dots

\dots nhưng lại là theo đúng nghĩa đen.

Nàng Pháp sư đang gặm một góc màn hình TV với bản mặt như một con thú dữ, còng nàng Cừu ôm lấy mép trên, mắt mơ màng như đang du hành vào cõi nào đó.

\img{10-4c}

``\vie{Vãi đạn!?}'' cậu Tổng bật ngửa, tay xòe ra như vừa chứng kiến một hiện tượng siêu nhiên.

``Hai đứa thậm chí còn chẳng xem kìa!'' nàng Y tá chỉ tay vào màn hình đang chiếu\dots\ phần giới thiệu chương trình phim truyện!

Kuon lập tức sắn tay áo, bước đến định kéo Rushia ra khỏi vô tuyến, trong khi Subaru bám lấy phần vai của nàng tóc xanh mà cố tách.

Nhưng hỡi ôi, Cô đồng vẫn bám chặt vào góc, như thể có một sức mạnh siêu nhiên nào đó khiến cho cô đơ cứng như đá vậy.

``Họ còn không xi nhê gì!'' Subaru thốt lên, hai tay tuyệt vọng cố kéo nàng Cô đồng khỏi chiếc TV.

``Em đã thử dùng chất tẩy rửa chưa?'' Choco đề xuất, nhặt lên một chai cồn lau tay sát bên.

Lập tức, cả hai người phản đối gay gắt. ``\textbf{Có phải là chất nhờn đâu!}'' nàng Vịt gào lên, mặt xanh như tàu lá.

``Thật. DÙng cái đó để mà chết người à!'' cậu Tổng ngán ngẩm đáp. ``Mà nếu dùng thì chắc gì hai đứa đã chịu buông?''

``Đúng đó! Chị xem này, cái con nhỏ này nó cắn luôn cả màn hình!!'' Subaru giật giật tay Rushia mà vẫn không nhúc nhích. ``Dính chặt luôn rồi!!''

Choco bình tĩnh nhìn xuống dưới bàn, rồi đáp với khuôn mặt giãn ra khi cô tìm được cách: ``Cô có cách rồi.''

Cô tiến đến gần, cầm điều khiển từ xa, ấn tắt nguồn.

*PHỤT*

Màn hình đen kịt. Ngay sau đó\dots

*UỴCH*

Rushia và Watame cùng lúc rớt phịch ra sau, ngã ngửa như hai khúc gỗ rụng. Nàng Pháp sư thậm chí còn run lên bần bật như bị chuột rút.

``Hai đứa là con thiêu thân à!?'' Subaru đứng chống nạnh, không thể tin vào mắt mình.

``Trông mấy đứa như mấy con nghiện vậy\dots'' Kuon đỡ trán, ngán ngẩm nhìn hai người họ đang lồm cồm bò dậy. ``Mà thôi, dậy đi. Nghe sensei giảng đạo kìa.''

Hai cô gái ngồi ủ rũ, co rúm lại trước nàng Y tá, đang đứng nghiêm trang như chuẩn bị viết giấy phạt.

``Mấy đứa đúng là hư thật. Cứ thế này thì chắc phải dỡ TV ra thôi.'' nàng Y tá lạnh lùng nói, khoanh tay nhìn về chiếc vô tuyến Subaru vừa bật lại.

``Tiếc tiền thì tiếc thật đấy,'' cậu Tổng chống hông chưng hửng tiếp lời, ``nhưng nói mãi không nghe thì phải làm vậy thôi.''

Vừa nghe đến từ ``dỡ TV'', cả hai cô gái bật dậy như lò xo, ánh mắt hoảng loạn. Họ lao tới hai người lớn, cầu xin thảm thiết.

\img{10-4d}

``\textbf{Không!! Đừng mà!!}'' Watame ôm lấy chân Kuon, đôi mắt xoay như chong chóng.

``Huốn hem hi-hi hơ! (Muốn xem ti-vi cơ!)'' Rushia mếu máo, tay nắm lấy vạt áo choàng của Choco.

Nàng Y tá và cậu Tổng nhìn xuống khuôn mặt lem nhem, ướt nước của hai đứa trẻ, trông như hai kẻ lang thang bị bắt quả tang ăn vụng.

``\textit{Trông mặt họ nhem nhuốc chưa kìa\dots}'' cả hai nghĩ thầm, đổ mồ hôi hột.

Sau một hồi đấu tranh nội tâm, hai người lớn đành phải xuống nước.

``Còn dính mặt lần nữa là khỏi ti-vi ti-vủng luôn nha,'' Choco giơ ngón trỏ, giọng nghiêm khắc.

``Anh cho hai đứa cơ hội lần cuối đấy. Không nhân nhượng nữa đâu,'' Kuon tiếp lời, ánh mắt cảnh cáo.

Khuôn mặt hân hoan của hai cô gái đã quay trở lại. Họ lập tức gật đầu như gà mổ thóc, vội vàng ngồi vào ghế với khoảng cách đúng chuẩn ba mét so với màn hình.

``Vâng ạ!''

``Bọn em sẽ ngồi xa ra!''

Kuon khoanh tay nhìn cảnh tượng trước mắt, thở môt hơi dài ngao ngán: ``\vie{Mùa thu đúng là mùa gặt drama.}''

\begin{center}
  \uline{Ngày hôm sau.}
\end{center}

Trời thu vẫn tiếp tục trải dài không gian bằng lớp sương mỏng nhè nhẹ trên các mái nhà. Lá đỏ rơi rụng dọc lối đi dẫn vào trụ sở, thỉnh thoảng lạo xạo dưới chân người đi qua.

Choco và Kuon lại cùng nhau đến văn phòng như thường lệ, mỗi người một tay cầm cốc cà phê nóng hổi. Đứng bên ngoài cửa, cả hai tạm dừng lại để trao đổi trước khi bước vào.

``Em vẫn thấy chuyện hôm qua khó tin thật\dots'' Choco khẽ nói, giọng nửa trách móc nửa buồn cười. ``Cắn TV thì em chưa thấy ai ngoài hai đứa đó.''

``Anh nghĩ đây là một trong những lần hiếm hoi mà não anh bị đơ khi thấy cảnh tượng kỳ quặc kia\dots''

``Có khi chúng ta phải cấm túc hai đứa thật.''

Kuon thở dài, quay sang hỏi: ``Em có chắc là hôm nay họ không dán mắt vào TV nữa không?''

``Để em ngó thử trước\dots'' Nàng Y tá khẽ hé cánh cửa phòng.

Bên trong, Rushia và Watame đang ngồi ngoan ngoãn trên hai chiếc ghế bành, lưng thẳng như học sinh tiểu học trong giờ đạo đức. Mắt không dán vào màn hình, tay đặt nghiêm chỉnh trên đùi. Tưởng như một bức tranh học đường.

``\dots Trông ổn áp phết nhỉ,'' cậu Tổng nhận xét, nét mặt thoáng bất ngờ.

Vừa mới dứt câu khen\dots

``\textbf{Hắt xì\dots}''

*COONG!* *COONG!*

Hai chậu nhôm từ trần rơi trúng đầu của cả Rushia và Watame, chẳng ai rõ sao chúng lại ở đó. Cả hai giật mình, ôm đầu mà không hề rời ghế.

\begin{figure}[H]
  \centering
  \begin{subfigure}[t]{0.45\textwidth}
    \centering
    \includegraphics[width=8cm]{../../../../Image/Book?/10-4e1.png}
  \end{subfigure}
  \quad
  \begin{subfigure}[t]{0.45\textwidth}
    \centering
    \includegraphics[width=8cm]{../../../../Image/Book?/10-4e2.png}
  \end{subfigure}
\end{figure}

``\hl{Nhìn như thể là rạp xiếc không bằng ấy\dots}'' Kuon đứng hình, ngơ ngác nhìn hai cô gái đang ôm đầu đau đớn.

``\hl{Em nghĩ chắc họ đang cố tỏ ra ngoan ngoãn\dots{} à mà thôi\dots}'' Choco lắc đầu. ``\hl{Hai đứa ấy định bàn nhau làm gì đây\dots}''

Bỗng, nàng Cừu liếc đồng hồ đeo tay, cầm chiếc điều khiển lên, rồi huých nhẹ nàng Pháp sư: ``Ru-chan, sắp đến giờ rồi đấy.''

Nàng Y tá cau mày, thầm nghĩ: ``\hl{Giờ xem TV à? Mà\dots{} sao họ trông mong chờ thế?}''

Và rồi, cả hai cô gái cùng nhau đồng thanh: ``\textbf{Đến giờ chương trình `Câu lạc bộ Lục lạc' rồi!!}''

Kuon nheo mắt, miệng lẩm bẩm: ``\hl{\dots Chẳng trách. Chuẩn cơm mẹ nấu rồi.}''

``\hl{Ý anh là sao vậy?}'' Choco hỏi.

``\hl{Hôm nọ anh thấy có cái lục lạc đặt trên bàn. Anh nghĩ nó phải dính dáng đến lý do vì sao họ mê TV thế chứ.}''

``\hl{Trời đất\dots{} Mà cái tên chương trình nghe ngớ ngẩn quá trời.}''

``\hl{Cái này thì phải hỏi nhà đài chứ đừng hỏi anh.}''

Trong phòng, chiếc TV bật sáng. Màn hình rực rỡ hiện lên hình ảnh hoạt hình dễ thương của một chú dê hồng đội nơ đỏ, tay cầm lục lạc, vừa hát vừa nhảy múa.

Và rồi hai cô gái cũng nhập vai theo.

Rushia và Watame bật khỏi ghế, lắc lư theo giai điệu, miệng ngân nga như trẻ mẫu giáo:

\begin{center}
  「ヘイ！へイへイ！タンバリンクラブ！
  あいつの明るさサンシャイン！」
  (Hây, hây hây! Câu lạc bộ lục lạc!
  Trời sáng chói, như là ban mai đây!)
\end{center}

\begin{figure}[H]
  \centering
  \begin{subfigure}[t]{0.45\textwidth}
    \centering
    \animategraphics[autoplay,loop,width=8cm]{12}{../../../../Image/Book?/10-4f/10-4f.}{1}{82}
    \caption{Họ vừa hát vừa gõ lục lạc loạn xạ vừa nhấp nhổm\dots}
  \end{subfigure}
  \quad
  \begin{subfigure}[t]{0.45\textwidth}
    \centering
    \animategraphics[autoplay,loop,width=8cm]{12}{../../../../Image/Book?/10-4f/10-4f.}{83}{191}
    \caption{\dots xoay người như cánh quạt\dots}
  \end{subfigure}
  \quad
  \begin{subfigure}[t]{0.45\textwidth}
    \centering
    \animategraphics[autoplay,loop,width=8cm]{12}{../../../../Image/Book?/10-4f/10-4f.}{192}{247}
    \caption{\dots đung đưa như những chú gà con\dots}
  \end{subfigure}
  \quad
  \begin{subfigure}[t]{0.45\textwidth}
    \centering
    \animategraphics[autoplay,loop,width=8cm]{12}{../../../../Image/Book?/10-4f/10-4f.}{248}{273}
    \caption{\dots và ưỡn người ra phía sau như nhện trong phim kinh dị, trong khi vẫn hát!}
  \end{subfigure}
\end{figure}

Choco-sensei lặng lẽ đóng cửa lại, lưng dựa vào cánh gỗ, ánh mắt như người vừa thấy ma. Trong khi đó, Kuon bịt miệng cố nhịn cười, nhưng thất bại thảm hại.

``\dots Em vừa xem cái gì vậy\dots?'' Choco hỏi, giọng run run. ``Trông họ\dots\ vui thật\dots''

``Anh nghĩ đó là lý do vì sao họ nghiện chương trình đó đấy. Tới mức dán mặt vào màn hình kia mà.''

``Chương trình con nít đó sao!?''

``Ừ,'' cậu gật gù. ``Anh mà là con nít anh còn thấy hay, huống gì hai đứa ấy.''

Cả hai cùng nhìn nhau, rồi lại nhìn cánh cửa đã đóng kín. Có vẻ như họ đã tìm ra lời giải cho hành vi kỳ quặc của Rushia và Watame trong mấy ngày gần đây.

Vấn đề là\dots

Họ sẽ xử trí như thế nào đây?

Vì cái thứ được gọi là ``Câu lạc bộ Lục lạc''\dots

\dots có lẽ sẽ trở thành nỗi ám ảnh mùa thu năm nay.
%==========================================TRUYỆN PHỤ=========================================
\omk

Từ bên ngoài, dù cửa đã khép, giọng hát trong trẻo của Rushia và Watame vẫn lọt ra, vang vang như gió thu đùa qua cửa sổ.

``\textbf{Hây! Hây hây! Câu lạc bộ Lục lạc! Trời sáng chói như là ban mai đây!}''

Kuon, ban đầu còn chỉ đứng dựa tường lắng nghe, bỗng gõ nhẹ nhịp lên tay, đung đưa vai theo tiết tấu, miệng vô thức hát lại: ``Hây\dots\ Hây hây\dots\ Câu lạc bộ lục lạc\dots''

Choco tròn mắt, nghiêng đầu nhìn cậu bằng ánh mắt kinh ngạc: ``Anh có vẻ cũng\dots\ thích thú nhỉ?''

``Lâu rồi mới quay về tuổi thơ một lần,'' cậu đáp, thoáng nở một nụ cười.

``Ể\dots?''

Cậu ngước lên trần nhà, nhớ lại những kỷ niệm hồi còn  thơ ấu: ``Giống như mấy bài hát mà các cô giáo mầm non hay hát cho học sinh ấy. Nghe xong tự nhiên thấy thư giãn thật\dots''

``Ý anh là\dots\ em nên vào trong đó múa phụ họa luôn à?''

``Nếu em muốn. Không phải sợ đâu. Đôi lúc quay về hồi còn bé con một lần cũng có chết ai đâu.''

Choco vừa định phản bác thì cậu con trai đã đẩy cửa vào, hiên ngang như đang bước lên sân khấu\dots\ và cảnh tượng bên trong khiến cậu suýt té ngửa.

Rushia đang bị Watame nâng chân lên cao, tư thế như đang ở dạng trồng cây chuối mà vẫn hát tiếp. Còn Watame thì một tay lắc lục lạc, một tay giữ chân đồng nghiệp, mắt vẫn hướng lên TV với sự kính cẩn như đang nhìn đấng linh thiêng.

``\textbf{A! Kuon-senpai!}'' nàng Cừu reo lên như vừa thấy một người khác muốn vào góp vui.

``Bọn em đang nhảy điệu Lục lạc đó!'' nàng Pháp sư hô vang, rồi chỉ về phía bàn. ``Anh muốn thử không? Bọn em có chuẩn bị thêm hai cái lục lạc nữa đó!''

Nàng Y tá vẫn thập thò ngoài cửa, chưa dám bước vào. Nhưng Watame không để cô trốn: ``Sensei ơi! Vào đây nhảy với bọn em đi!''

``Ơ\dots\ ơ\dots\ Được rồi\dots'' cô giáo lúng túng, nhưng rồi cậu Tổng ném lục lạc cho cô, cô bắt lấy như thể bản năng giáo viên đã thức dậy.

``\textbf{Hây, hây hây! Câu lạc bộ lục lạc\~{}!}'' Kuon bắt nhịp hô, khiến cả phòng như bừng lên sinh khí.

``Ê\dots\ liệu có ổn không vậy\dots?'' Choco lúng túng hỏi.

``Cô cứ nhảy theo bọn em thôi!'' Watame cười hớn hở.

``Anh nghĩ em nên làm chỉ huy luôn. Biết đâu lại vui đấy,'' cậu nháy mắt.

``\dots Thôi được rồi!''

Choco giơ cao lục lạc, hít sâu, rồi mạnh dạn hô: ``\textbf{Mấy đứa! Làm theo cô này!}''

Cả bốn người cùng nhảy múa, theo một điệu nhạc vốn dành cho trẻ mẫu giáo. Nhưng không ai trong số họ thấy kỳ cục cả. Trong khoảnh khắc ấy, họ chẳng còn là người lớn nữa, họ chỉ là những tâm hồn đang vui đùa giữa ánh nắng tháng Chín, nơi tuổi thơ và sự hồn nhiên chưa bị thế giới vấy bẩn.

Nàng Y tá ngắm nhìn gương mặt ngây ngô và rạng rỡ của Rushia và Watame, rồi quay sang Kuon đang gõ lục lạc theo kiểu nhịp trưởng nhà trẻ. Cô mỉm cười nhẹ: ``Thì ra là vậy. Thì ra vì thế mà mấy đứa say mê chương trình này\dots''

Cô không quên lên giọng nghiêm: ``Cơ mà vẫn phải để mắt ra xa TV đấy nhé. Không hỏng mắt là không được xem nữa đâu.''

``Vâng\~{}!'' hai cô gái đồng thanh.

``Anh chưa nghe rõ!'' Kuon hào hứng hô lớn, ``\textbf{Nhớ chưa!?}''

``\textbf{Vâng ạ!!}'' họ hét lên, tràn đầy sinh khí.

Họ nhảy thêm một lượt, rồi một lượt nữa. Cảm giác như thời gian cũng bị cuốn theo điệu lục lạc.

Cho đến khi\dots

``Này mấy người\dots\ Mọi người đang làm cái quái gì vậy\dots?''

Cánh cửa sau lưng họ bật mở, và Subaru xuất hiện với bản mặt ngơ ngác, tay cầm ổ bánh mì vẫn đang ăn dở. ``Đây có phải là tiết học giáo dục công dân không vậy, shuba!?''

``Subaru-chan/-senpai!?'' cả nhóm đồng thanh.

``Cái này là cái gì? `Lục lạc' là thứ gì vậy? Sao nhìn ai cũng như bị\dots\ thôi miên thế!?'' nàng Vịt vẫn không khỏi ngỡ ngàng, khi hai người lớn cuối cùng cũng bị ``chinh phục'' bởi điệu nhảy trẻ con.

Cậu Tổng vẫy tay, cười rạng rỡ: ``Vào đi. Em cũng cần một vé trở về tuổi thơ.''

Subaru nhìn họ thêm vài giây. Và rồi\dots\ đặt bánh mì xuống bàn. ``\dots Thôi kệ. Cho em một cái lục lạc.''

Và cứ thế, chương trình ``Câu lạc bộ Lục lạc'' có thêm một thành viên nữa.

Họ không biết đến bao giờ cơn sốt này mới hết.

Nhưng vào khoảnh khắc đó, họ chỉ biết một điều: ``Lắc lục lạc là vui. Và đôi khi\dots\ họ cũng có quyền quay trở về tuổi thơ, dù chỉ là chốc lát\dots''

\end{document}